% Motivation and learning outcomes.
% The outcomes are the outcomes of Day 13 in the syllabus.
\begin{frame}{Why this day matters}
  \begin{itemize}
    \item On Days 8 to 12 you built applications that detect objects, stream
          video, and send results with MQTT. A deployed application runs for
          weeks with no person in front of it.
    \item A deployed model can fail with no error message. The light changes,
          a new object appears, a sensor gets dirty, or the board gets hot.
          The program still runs, but its answers get worse.
    \item Monitoring makes such failures visible. The device measures itself,
          writes logs that fit its storage, and sends metrics to a dashboard.
          An alert tells a person before the users see the problem.
    \item In the lab you add measurements and logs to the Day 8 detector, show
          them on a live dashboard, cause a drift event, detect it from the
          confidence values, and test one alert rule.
  \end{itemize}
\end{frame}

\begin{frame}{Learning outcomes}
  After this day you can:
  \begin{itemize}
    \item Select metrics for model health and for system health.
    \item Write structured logs on a device with small storage.
    \item Build a live dashboard for an edge device.
    \item Detect data drift without labels.
  \end{itemize}
\end{frame}
